% ---- /md-to-pdf 共通プリアンブル（pLaTeX 用）----
% 学会指定クラス（ipsj.cls など）は pLaTeX 前提のものが多く、LuaLaTeX 用の
% md-common.tex は使えない（luatexja-adjust が読めない）。こちらを使う。
% レイアウト側で  %%! engine = platex  と  %%! common = md-common-platex.tex  を宣言する。

\usepackage{amsmath,amssymb}   % pandoc は別行立て数式に aligned を使う
\usepackage{graphicx}
\usepackage{enumitem}
\usepackage{booktabs,array}
\usepackage{calc}
\usepackage{url}

% 学会クラスは2段組が多く、longtable は1段組でしか使えない。pandoc は表を
% longtable で出すので、tabular に読み替える（ページをまたぐ長い表は組めない。
% その場合は .tex を手で直して table* で全幅に逃がす）。
\makeatletter
% \end... という名前は \providecommand で定義できないので \def を使う
\def\endhead{}
\def\endfirsthead{}
\def\endlastfoot{}
\def\endfoot{}
\newenvironment{longtable}[2][]{\par\centering\begin{tabular}{#2}}{\end{tabular}\par}
\makeatother

% 画像の相対パスを入力mdのフォルダ基準で解決する
\graphicspath{__GRAPHICSPATH__}

% pandoc は \includegraphics に alt= キーを付けるので graphicx に無視させる
\makeatletter
\define@key{Gin}{alt}[]{}
\makeatother

% ipsj.cls は \newblock を定義しない（openbib オプション時にだけ再定義する）。
% 一方 ipsjunsrt.bst は note フィールドを \newblock で区切って書き出すため、
% .bib に note があると「Undefined control sequence」でビルドが止まる。
% 標準クラスと同じ定義を補っておく。
\makeatletter
\providecommand{\newblock}{\hskip .11em\@plus.33em\@minus.07em}
\makeatother

% pandoc ヘルパ
\providecommand{\tightlist}{\setlength{\itemsep}{0pt}\setlength{\parskip}{0pt}}
\providecommand{\pandocbounded}[1]{\begin{center}#1\end{center}}
\providecommand{\hl}[1]{\textbf{#1}}   % ==ハイライト==（学会原稿では色を付けない）
\providecommand{\passthrough}[1]{#1}

% 箇条書きは深いネストもあり得るので上限を広げる
\setlistdepth{9}
\renewlist{itemize}{itemize}{9}
\renewlist{enumerate}{enumerate}{9}
% \renewlist はラベル定義も消すので、各階層に必ず設定し直す
\setlist[itemize]{label=\textbullet}
\setlist[itemize,2]{label=\normalfont\bfseries--}
\setlist[itemize,3]{label=\textasteriskcentered}
\setlist[itemize,4]{label=\textperiodcentered}
\setlist[itemize,5]{label=\textperiodcentered}
\setlist[itemize,6]{label=\textperiodcentered}
\setlist[itemize,7]{label=\textperiodcentered}
\setlist[itemize,8]{label=\textperiodcentered}
\setlist[itemize,9]{label=\textperiodcentered}
\setlist[enumerate]{label=(\arabic*)}
\setlist[enumerate,2]{label=(\alph*)}
\setlist[enumerate,3]{label=(\roman*)}
\setlist[enumerate,4]{label=(\arabic*)}
\setlist[enumerate,5]{label=(\alph*)}
\setlist[enumerate,6]{label=(\roman*)}
\setlist[enumerate,7]{label=(\arabic*)}
\setlist[enumerate,8]{label=(\alph*)}
\setlist[enumerate,9]{label=(\roman*)}

% pandoc --natbib は \citep / \citet を出すが、学会配布の .bst（ipsjunsrt 等）は
% natbib 非対応のことが多い。素の \cite に読み替える。
% 注意: ページ指定 [@key, 10--15] の頁番号はここで落ちる（学会の引用形式では
% 通常ページを示さないため）。示したい場合は本文に直接書く。
\makeatletter
\newcommand{\citep@fallback}[2][]{\cite{#2}}
\@ifundefined{citep}{\let\citep\citep@fallback}{}
\@ifundefined{citet}{\let\citet\citep@fallback}{}
\makeatother

% citeproc 引用（%%! citations = citeproc のとき）。natbib を使う場合は不要だが、
% 定義しておいても害はない。
\providecommand{\citeproctext}{}
\providecommand{\citeproc}[2]{\cite{#1}}
\newlength{\cslhangindent}\setlength{\cslhangindent}{1.5em}
\newlength{\csllabelwidth}\setlength{\csllabelwidth}{3em}
\newenvironment{CSLReferences}[2]
 {\begin{list}{}{%
   \setlength{\itemindent}{0pt}\setlength{\leftmargin}{0pt}\setlength{\parsep}{0pt}%
   \ifodd #1\setlength{\leftmargin}{\cslhangindent}\setlength{\itemindent}{-1\cslhangindent}\fi
   \setlength{\itemsep}{0.2\baselineskip}}}
 {\end{list}}
\providecommand{\CSLBlock}[1]{\hfill\break\parbox[t]{\linewidth}{\strut\ignorespaces#1\strut}}
\providecommand{\CSLLeftMargin}[1]{\parbox[t]{\csllabelwidth}{\strut#1\strut}}
\providecommand{\CSLRightInline}[1]{\parbox[t]{\linewidth}{\strut#1\strut}}
\providecommand{\CSLIndent}[1]{\hspace{\cslhangindent}#1}
